% Context from chapters-tex/fourier.tex; for mathematical reference, not compilation.
\section{Fourier in Multiple Dimensions}
\label{sec-fourier-multiple-dim}

Tensor products extend the one-dimensional Fourier transform to any finite dimension $d>1$.


                                                
\subsection{On Continuous Domains}

   
\paragraph{On \texorpdfstring{$\RR^d$}{Euclidean Space}.}

\begin{figure}[tbp]
\centering
\BookDrawing{tikz/fourier-wave}
\caption{2-D sine wave.}
\label{fig-review-fourier-wave}
\end{figure}
Tensor products retain a useful feature of Fourier atoms: a product of elementary 1-D complex exponentials is a plane wave
\eq{
	\prod_{\ell=1}^d e^{ \imath x_\ell \om_\ell } = e^{ \imath \dotp{x}{\om} }
}
whose constant-phase hyperplanes are orthogonal to the wave vector $\om=(\om_\ell)_{\ell=1}^d \in \RR^d$ when this vector is nonzero. Here $\dotp{x}{\om} = \sum_\ell x_\ell \om_\ell$ is the canonical inner product on $\RR^d$.

The Fourier transform and its inverse are defined by
\begin{align*}
	\foralls \om \in \RR^d, \quad \hat f(\om) &\eqdef \int_{\RR^d} f(x) e^{-\imath\dotp{x}{\om}} \d x, \\
	\foralls x \in \RR^d, \quad  f(x) &= \frac{1}{(2\pi)^d} \int_{\RR^d}\hat f(\om)e^{\imath\dotp{x}{\om}}\,\d\om,
\end{align*}
under the same integrability assumptions as in one dimension.


\myfigure{
\image{orthobases}{.3}{2d-extension-fourier-plane}
\image{orthobases}{.2}{2d-extension-fourier-1}
\image{orthobases}{.2}{2d-extension-fourier-2}
\image{orthobases}{.2}{2d-extension-fourier-3}
}{ 
	2D Fourier orthogonal bases.   
}{fig-2d-fourier-extension}


   
\paragraph{On \texorpdfstring{$(\RR/2\pi\ZZ)^d$}{the Torus}.}

Given an orthonormal basis $(\phi_{n_1})_{n_1 \in \NN}$ of $L^2(\XX)$, one constructs an orthonormal basis of $L^2(\XX^d)$ by tensorization
\eql{\label{eq-tensor-product}
	\foralls k=(k_1,\ldots,k_d) \in \NN^d, \quad
	\foralls x \in \XX^d, \quad
	\phi_k(x) = \phi_{k_1}(x_1) \ldots \phi_{k_d}(x_d).
}
Fubini's theorem gives orthogonality. Completeness gives convergence of the rectangular partial sums $\sum_{\norm{k}_\infty \leq N} \dotp{f}{\phi_k} \phi_k \rightarrow f$ in $L^2(\XX^d)$.

           
\begin{figure}[tbp]
\centering
\BookDrawing{tikz/fourier-torus}
\caption{The two-dimensional torus $\TT^2=(\RR/2\pi\ZZ)^2$.}
\label{fig-review-fourier-torus}
\end{figure}
   
For the multidimensional torus $(\RR/2\pi\ZZ)^d$, using the Fourier basis~\eqref{eq-fourier-basis-1d}, this gives the basis
\eq{
	\foralls k \in \ZZ^d, \quad \phi_k(x) = e^{\imath \dotp{x}{k}}
} 
which is indeed a Hilbertian orthonormal basis for the inner product $\dotp{f}{g} \eqdef \frac{1}{(2\pi)^d} \int_{\TT^d} f(x)\bar g(x)\d x$. 
 
This defines the Fourier transform and the reconstruction formula on $L^2(\TT^d)$
\eq{
	\hat f_k \eqdef \frac{1}{(2\pi)^d} \int_{\TT^d} f(x) e^{-\imath \dotp{x}{k}}  \d x
	\qandq
	f = \sum_{k \in \ZZ^d} \hat f_k e^{\imath \dotp{x}{k}}.
}

                                                
\subsection{On Discrete Domains}
\label{sec-fft-2d}

    
\paragraph{Discrete Fourier Transform.}

\begin{figure}[tbp]
\centering
\BookDrawing{tikz/fourier-torus-discrete}
\caption{Discrete 2-D torus.}
\label{fig-review-fourier-torus-discrete}
\end{figure}
On a $d$-dimensional discrete domain of the form
\eq{
	n = (n_1,\ldots,n_d) \in \YY_d \eqdef \range{1,N_1} \times \ldots \times \range{1,N_d}
}
(we denote $\range{a,b} \eqdef \enscond{i \in \ZZ}{a \leq i \leq b}$)
of $N=N_1\ldots N_d$ points, with periodic boundary conditions, one defines an orthogonal basis $(\phi_k)_k$ by the same tensor product formula as~\eqref{eq-tensor-product} but using the 1-D discrete Fourier basis~\eqref{eq-dft}
\eql{\label{eq-exp-four-disc-multid}
	\foralls (k,n) \in \YY_d^2, \quad
	\phi_k(n) = \phi_{k_1}(n_1) \ldots \phi_{k_d}(n_d) 
	= \prod_{\ell=1}^d e^{\frac{2\imath\pi}{N_\ell} k_\ell n_\ell}
	= e^{ 2\imath\pi \dotp{k}{n}_{\YY_d} } 	
} 
where we used the (rescaled) inner product
\eql{\label{eq-inner-prod-multi-d}
	\dotp{k}{n}_{\YY_d} \eqdef \sum_{\ell=1}^d \frac{k_\ell n_\ell}{N_\ell}.
}
The pairing in~\eqref{eq-inner-prod-multi-d} acts on frequency and spatial index vectors. The basis $(\phi_k)_k$ is orthonormal for the signal inner product $N^{-1}\sum_n f_n\bar g_n$. As in one dimension, we define the Fourier coefficients without normalizing by $(N_\ell)_\ell$